\documentclass[../../main.tex]{subfiles}
\begin{document}

\begin{flushright}
\footnotesize\textit{【自動文字起こし・要確認】}
\end{flushright}

{\bf [解]}
\begin{align}
\begin{cases}
b_1 = 2 \\
b_{n+1} = b_n + \dfrac{1}{b_n} + 2
\end{cases} \quad \cdots \text{①}
\end{align}

\paragraph{(1)}

$b_n > 2n \quad (n \ge 2) \quad \cdots \Diamond$ を帰納法で示す．$n=2$ の時，①から $b_2 = 4 + \dfrac{1}{2}$ なので $\Diamond$ は成立．$n=k \in \mathbb{N}_{\ge 2}$ での成立をカテイすると，①から
\begin{align}
b_{k+1} = b_k + \dfrac{1}{b_k} + 2 > 2(k + 1)
\end{align}
となり，$n=k+1$ でも $\Diamond$ は成立．よって示された．(終)

\paragraph{(2)}

(1) から $n \ge 2$ に対して $\dfrac{1}{a_n} > 2n \ \therefore \ a_n < \dfrac{1}{2n}$ である．$n \to \infty$ をかんがえるので $n$ は十分大きいとして良く，右図で面積比較して，
\begin{align}
\dfrac{1}{n} \sum_{k=1}^n a_k &< \dfrac{1}{n} \left( \dfrac{1}{2} + \sum_{k=2}^n \dfrac{1}{2k} \right) \\
&< \dfrac{1}{2n} \left( 1 + \int_1^n \dfrac{1}{x} dx \right) \\
\therefore \ \dfrac{1}{n} \sum_{k=1}^n a_k &< \dfrac{1 + \log n}{2n} \to 0 \quad (n \to \infty)
\end{align}
から，$0 < a_k$ とあわせて ($\because$ ①から帰納的に $0 < a_k$)，はさみうちから，
\begin{align}
\dfrac{1}{n} \sum_{k=1}^n a_k \to 0 \quad (\text{終})
\end{align}

\begin{figure}[htb]
\centering
\begin{tikzpicture}
\draw[->] (0,0) -- (3,0) node[right]{$x$};
\draw[->] (0,0) -- (0,2) node[above]{$y$};
\draw[domain=0.5:2.5, smooth] plot (\x, {1/\x});
\draw (1,0) rectangle (1.5, 1);
\draw (1.5,0) rectangle (2, 0.66);
\node[below] at (1,0) {$1$};
\node[below] at (2,0) {$n$};
\node at (0,1) [left] {$\frac{1}{a_n}$}; % rough inclusion based on scratch
\end{tikzpicture}
\caption{$y=\dfrac{1}{x}$ のグラフと区分求積による面積の評価}
\end{figure}

\paragraph{(3)}

①から，
\begin{align}
b_{n+1} - b_n = 2 + a_n
\end{align}
$n=1, 2, \cdots, N$ として足して，
\begin{align}
b_{N+1} - b_1 = 2N + \sum_{k=1}^N a_k
\end{align}
両辺 $N+1$ でわって，
\begin{align}
\dfrac{1}{(N+1)a_{N+1}} &= \dfrac{2}{N+1} + 2\dfrac{N}{N+1} \cdot \dfrac{1}{N} \sum_{k=1}^N a_k \\
&\xrightarrow{N \to \infty} 2 \quad (\because (2))
\end{align}
だから，$0 < x$ で $\dfrac{1}{x}$ が連続なことから，
\begin{align}
N a_N \to \dfrac{1}{2} \quad (N \to \infty) \quad (\text{終})
\end{align}

\bigskip
\noindent{\bf [解2]}
\paragraph{(3)}

$b_n \le 2n + \log n$ を帰納法で示す．$n=1$ の時は成立するので，$1 \le n = k \in \mathbb{N}$ での成立をカテイ．①から
\begin{align}
b_{k+1} = b_k + \dfrac{1}{b_k} + 2 \le 2(k+1) + \log k + \dfrac{1}{2k} \quad (\because 2k \le b_k) \quad \cdots \text{②}
\end{align}
ここで，$f(x) = \log(x+1) - \log x - \dfrac{1}{2x}$ とすると，
\begin{align}
f'(x) = \dfrac{1}{x+1} - \dfrac{1}{x} + \dfrac{1}{2x^2} = \dfrac{1-x}{2x^2(x+1)}
\end{align}
から増減表をうる
\begin{table}[htb]
\centering
\begin{tabular}{|c|c|c|c|c|}
\hline
$x$ & $0$ & $\cdots$ & $1$ & $\cdots$ \\ \hline
$f'(x)$ & & $+$ & $0$ & $-$ \\ \hline
$f(x)$ & & $\nearrow$ & & $\searrow$ \\ \hline
\end{tabular}
\caption{$f(x)$ の増減表}
\end{table}
$\left( \text{$x \to \infty$ の時，平均値から $x < c < x+1$ なる $c$ に対し，} f(x) = \dfrac{1}{c} - \dfrac{1}{2x} > \dfrac{1}{x+1} - \dfrac{1}{2x} > 0 \right)$

よって $1 \le x$ では $f(x) > 0$ つまり $\log(k+1) > \log k + \dfrac{1}{2k}$ だから，②より
\begin{align}
b_{k+1} \le 2(k+1) + \log(k+1)
\end{align}
よって $n=k+1$ でも成立．以上から成立し，
\begin{align}
2n < \dfrac{1}{a_n} < 2n + \log n
\end{align}
だから，はさみうちより $n a_n \to \dfrac{1}{2}$ (終)

\end{document}
